\clearpage\pagestyle{otMain}\fontsize{11}{15}\selectfont\corrigesChapitre{9}{Fonctions linéaires et affines}
\begin{blocLivre}{cor-9-1}\begin{corrige}{1}{Reconnaître les formes}\demarche{Comparer les expressions aux définitions.}
Les fonctions $\res{f,g,k}$ sont affines, car $k(x)=0x+7$. Seule $\res f$ est linéaire parmi elles. La fonction carré $h$ n’est pas affine.
\end{corrige}\end{blocLivre}
\begin{blocLivre}{cor-9-2}\begin{corrige}{2}{Calculer avec une linéaire}\demarche{Appliquer puis inverser la multiplication.}
$f(4)=-2{,}5\times4=\res{-10}$. L’égalité $-2{,}5x=15$ donne $\res{x=-6}$. On divise par un coefficient non nul.
\end{corrige}\end{blocLivre}
\begin{blocLivre}{cor-9-3}\begin{corrige}{3}{Lire les coefficients}\demarche{Relier les coefficients à leur rôle graphique.}
Le coefficient directeur est $\res{-3}$, l’ordonnée à l’origine $\res8$ et $g(2)=-6+8=\res2$. La droite descend de gauche à droite.
\end{corrige}\end{blocLivre}
\begin{blocLivre}{cor-9-4}\begin{corrige}{4}{Tracer une linéaire}\demarche{Choisir une entrée donnant une image simple.}
La droite passe par l’origine $O(0\,;0)$. Une autre image facile est $f(2)=3$, qui donne $A(2\,;3)$. On trace la droite passant par ces deux points ; le corrigé la montre en vert.
\label{fig-cor-9-4}\figureGeo[corrige]{affine-trace-lineaire}
\end{corrige}\end{blocLivre}
\begin{blocLivre}{cor-9-5}\begin{corrige}{5}{Tracer une affine}\demarche{Prendre les intersections avec les axes quand elles sont simples.}
On peut utiliser $(0\,;4)$ et $(4\,;0)$. L’antécédent de 0 est $\res4$, puisque $-4+4=0$. Le graphique et l’équation concordent.
\label{fig-cor-9-5}\figureGeo[corrige]{affine-trace}
\end{corrige}\end{blocLivre}
\begin{blocLivre}{cor-9-6}\begin{corrige}{6}{Une linéaire à retrouver}\demarche{Utiliser la forme annoncée avant les données numériques.}
On écrit $f(x)=ax$. Comme $6a=21$, $a=7/2$. Donc $\res{f(x)=3{,}5x}$ et $\res{f(-2)=-7}$. L’information « linéaire » est essentielle pour utiliser un seul point.
\end{corrige}\end{blocLivre}
\begin{blocLivre}{cor-9-7}\begin{corrige}{7}{Deux points, une formule}\demarche{Calculer le coefficient directeur puis utiliser un point pour la constante.}
Le coefficient directeur est $(10-1)/(5-2)=3$. On remplace les coordonnées de A dans $y=3x+b$ : $1=6+b$, donc $b=-5$. Ainsi $\res{f(x)=3x-5}$. On contrôle $f(5)=10$.
\label{fig-cor-9-7}\figureGeo[corrige]{affine-deux-points}
\end{corrige}\end{blocLivre}
\begin{blocLivre}{cor-9-8}\begin{corrige}{8}{Retrouver un tarif}\demarche{Interpréter la variation de prix avant de retrouver le forfait.}
L’augmentation de 12 € correspond à trois heures : le prix horaire est $\res{4\text{ €}}$. Pour deux heures, le forfait vaut $14-2\times4=\res{6\text{ €}}$. Le modèle est $P(h)=4h+6$.
\end{corrige}\end{blocLivre}
\begin{blocLivre}{cor-9-9}\begin{corrige}{9}{Coefficients lus sur un graphique}\demarche{Lire l’ordonnée à l’origine, puis la variation pour une unité d’abscisse.}
La droite coupe l’axe des ordonnées en 3 : $\res{b=3}$. Quand l’abscisse augmente de 1, par exemple de 0 à 1, l’ordonnée passe de 3 à 1 : elle diminue de 2, donc $\res{a=-2}$. La fonction est $\res{f(x)=-2x+3}$ ; la droite descend, ce qui confirme $a<0$.
\label{fig-cor-9-9}\figureGeo[corrige]{affine-lecture}
\end{corrige}\end{blocLivre}
\begin{blocLivre}{cor-9-10}\begin{corrige}{10}{Comparer par une différence}\demarche{Calculer une différence indépendante de l’entrée.}
La différence $g(x)-f(x)=7$ est toujours positive. Donc $\res{g(x)=f(x)+7}$ pour tout $x$. Les droites ont le même coefficient directeur et des ordonnées à l’origine différentes : elles sont parallèles.
\label{fig-cor-9-10}\figureGeo[corrige]{affine-paralleles}
\end{corrige}\end{blocLivre}
\begin{blocLivre}{cor-9-11}\begin{corrige}{11}{Deux abonnements}\demarche{Résoudre le seuil, puis l’inéquation, et revenir aux nombres entiers.}
L’égalité $20+3n=8+5n$ donne $\res{n=6}$ et un coût de 38 €. L’inéquation équivaut à $12<2n$, soit $n>6$ : A est moins chère $\res{\text{à partir de 7 séances}}$, B l’est pour moins de 6 séances. Les droites se croisent au point $(6\,;38)$.
\label{fig-cor-9-11}\figureGeo[corrige]{affine-tarifs}
\end{corrige}\end{blocLivre}
\begin{blocLivre}{cor-9-12}\begin{corrige}{12}{Une vidange à débit constant}\demarche{Associer la formule à son domaine physique.}
Le modèle est $\res{V(t)=900-35t}$. L’égalité $900-35t=200$ donne $\res{t=20\text{ min}}$. À 30 minutes, la formule donne $-150$ L, impossible : le modèle s’arrête quand la cuve est vide, vers 25,7 minutes.
\end{corrige}\end{blocLivre}
\begin{blocLivre}{cor-9-13}\begin{corrige}{13}{Une vitesse et une inéquation}\demarche{Modéliser par une fonction linéaire, puis résoudre l’inéquation.}
La distance est proportionnelle à la durée : $\res{d(t)=5t}$, fonction linéaire. L’inéquation $5t>120$ équivaut à $\res{t>24}$, car 5 est positif. Graphiquement, la droite $y=5t$ passe au-dessus de la droite horizontale $y=120$ après l’abscisse 24.
\end{corrige}\end{blocLivre}
\begin{blocLivre}{cor-9-14}\begin{corrige}{14}{Un point suffit-il ?}\demarche{Compter les informations et construire deux exemples.}
Non : $f(x)=x+5$ et $g(x)=3x+1$ donnent toutes deux 7 en 2, mais diffèrent pour d’autres entrées. Une condition unique ne détermine pas les deux coefficients. Pour une fonction linéaire, $b$ serait déjà imposé à zéro.
\end{corrige}\end{blocLivre}
\begin{blocLivre}{cor-9-15}\begin{corrige}{15}{Un seuil avec contrainte}\demarche{Résoudre chaque inéquation, puis prendre le plus grand entier solution.}
À 6 jours, A vaut 117 € et B 123 € : seule A respecte le budget. Pour A, $45+12d\le120$ donne $d\le6{,}25$ : le maximum est $\res6$ jours. Pour B, $15+18d\le120$ donne $d\le5{,}83\ldots$ : le maximum est $\res5$ jours. On ne garde que les durées entières.
\end{corrige}\end{blocLivre}
